\documentclass[14pt]{extarticle}
\usepackage{preamble}
\geometry{letterpaper, landscape, margin=0.5in}
\renewcommand{\answer}[1]{}
\setlength{\parskip}{4pt}

\begin{document}
\pagestyle{empty}

\daybanner{Unit 3 \textperiodcentered\ Topic 3.14 \textperiodcentered\ THE PROBLEM}{One Table, Two Study Designs}

\noindent\begin{minipage}[t]{0.57\linewidth}
\vspace{0pt}
Suppose a transit authority obtains the shown counts under either design.\par\smallskip \textbf{Design A:} Take independent SRSs of 100 from each of three route populations of 5,000 and record ticket method.\par \textbf{Design B:} Take one SRS of 300 from 15,000 system riders and record both route and ticket method.\par\smallskip \textbf{Ari:} ``Route rows imply homogeneity; the counts are enough to choose.''\par \textbf{Bea:} ``Read how riders were sampled; identical counts can support different population claims.''\par
\end{minipage}\hfill
\begin{minipage}[t]{0.40\linewidth}
\vspace{0pt}
\begin{center}
\textbf{Ticket method}\par\smallskip
\begin{tabular}{|l|c|c|c|c|}
\hline
\textbf{Route} & \textbf{Mobile} & \textbf{Card} & \textbf{Cash} & \textbf{Total} \\ \hline
\textbf{North} & 60 & 25 & 15 & 100 \\ \hline
\textbf{East} & 48 & 32 & 20 & 100 \\ \hline
\textbf{West} & 42 & 33 & 25 & 100 \\ \hline
\textbf{Total} & 150 & 90 & 60 & 300 \\ \hline
\end{tabular}
\end{center}
\end{minipage}

\begin{firsttakebox}{FIRST TAKE --- one sentence before you work the parts}
Do the same counts mean that riders were selected in the same way? Explain in one sentence.
\end{firsttakebox}

\begin{description}[leftmargin=1.15in, labelwidth=1in, itemsep=2pt, topsep=2pt]
  \item[\dokbadge{1}~(a)] How many random samples are taken in Design A and in Design B?
  \item[\dokbadge{2}~(b)] Calculate the three expected counts for the North row. Explain briefly why the other two rows have the same expected counts.
  \item[\dokbadge{3}~(c)] In 4 sentences, judge Ari's rule. Name and justify the test for each design, state each null hypothesis in context, and explain why the same table can lead to different test names.
\end{description}

\vfill
\textbf{Self-paced bonus problem. Work (a)--(c) from this sheet; turn it in whenever you finish.}

\end{document}
